% Standalone section fragment. Requires amsmath, amssymb, amsthm.
% The parent report should define theorem/proposition/proof environments.
\section{A variable-alias obstruction to automatic freshness padding}
\label{sec:padding-obstruction}

The following result rules out a specific tempting proof of guard
elimination. It concerns derivability of the old gate in the proposed
modified numerical calculus. It does not prove that every guardless
calculus has insufficient strength for a PA-consistency construction.

\subsection{Exact calculus and the proposed replacement}

Put
\[
 p(a,b)=\frac{(a+b)(a+b+1)}2+b,\qquad
 I(A,B)=4p(A,B),\qquad M(a,b)=4p(a,b)+1,
\]
\[
 U(a,B)=4p(B,a)+2,\qquad E(a,b)=4p(a,b)+3.
\]
Code $0$ is truth. A nonzero code is decoded by its remainder modulo
$4$ into implication, membership, universal quantification, or equality,
respectively. In particular $1=M(0,0)$ is the designated false formula,
and $I(0,0)=0$ is a code alias. Write $N(A)=I(A,1)$; conjunction and
existential quantification are the usual classical abbreviations formed
from $I,N,U$. In displayed formulas we use logical connective notation
for these numerical constructors.

The unchanged heads are the following schemata, for arbitrary numerical
parameters of the indicated sorts:
\begin{align*}
 &(\exists x_a\,\forall x_a C)\to C, \tag{H1}\\
 &U(a,U(b,C))\to U(b,U(a,C)), \tag{H2}\\
 &M(c,a)\to(E(c,b)\to M(b,a)), \tag{H3}\\
 &E(c,a)\to(M(b,c)\to M(b,a)), \tag{H4}\\
 &\exists x_a E(c,a), \tag{H5}\\
 &U(a,B)\to\bigl(U(a,B\to C)\to U(a,C)\bigr). \tag{H6}
\end{align*}
Also retain truth, the dominated propositional head
\[
 ((H\to d)\to D)\to((D\to H)\to(A\to H)),\qquad d\le D,
\]
and the unrestricted separate diagonal head
\[
 \Gamma(B)=B\to U(B,B).
\]
The inference rules are modus ponens and generalization of theorems.
No additional rule allowing arbitrary renaming of numerical variable
labels is postulated.

Define the old gate by
\begin{align*}
 V(a,b,r)&=M(a,b)\to U(r,M(a,b)),\\
 P(a,b,r)&=r\to M(r,b),\\
 S(a,b,r)&=U\bigl(b,
    (M(a,b)\to P(a,b,r))\wedge(M(r,b)\to r)\bigr),\\
 Q(a,b,r)&=U(r,S(a,b,r)\to N(P(a,b,r))),\\
 G(a,b,r)&=N(V(a,b,r)\to Q(a,b,r)).
\end{align*}
Occurrences of $r$ as a formula mean the formula decoded from code $r$;
its occurrences as an argument of $M$ or the first argument of $U$ are
variable labels. The original gate schema supplies $G(a,b,r)$ whenever
$r\ne a,b$.

The proposed unguarded replacement uses
\[
 F(a,b,c)=I\bigl(I(M(a,b),M(a,b)),c\bigr)
\]
and supplies only
\[
 G\bigl(a,b,F(a,b,c)\bigr)
 \quad\text{for all }a,b,c\in\mathbb N.
 \tag{PG}
\]
Let $\mathcal C_{\mathrm{pad}}$ be the resulting calculus, with all the
unchanged heads above and (PG).

The numerical freshness calculation is correct. If $m=M(a,b)$, then
$m>a,b$ and $m\ge1$. Consequently $I(m,m)\ge16$, and
\[
 F(a,b,c)>a,b,c,\qquad F(a,b,c)\ge F(0,0,0)=544.
 \tag{1}
\]
Moreover $I(m,m)$ is a propositional theorem, so the formula with code
$F(a,b,c)$ is propositionally equivalent to the formula with code $c$.
Nevertheless, these facts do not recover the old gate schema.

\begin{theorem}[Failure of old-gate recovery]
\label{thm:padding-alias}
In the calculus $\mathcal C_{\mathrm{pad}}$, neither
\[
 V(1,0,4)=M(1,0)\to U(4,M(1,0))
\]
nor the legitimate original gate instance $G(1,0,4)$ is derivable.
In particular, unrestricted separate diagonal axioms do not repair the
loss of arbitrary-label membership vacuity caused by this padding scheme.
\end{theorem}

\begin{proof}
We give a sound interpretation of every axiom and rule of
$\mathcal C_{\mathrm{pad}}$ in which the displayed formulas fail.

Let the object domain be $\mathbb N$. Write $\mu(s,z)$ when the $z$-th
binary digit of $s$ is $1$. Thus $\mu(s,z)$ means that the finite set
coded by $s$ contains $z$. In particular
\[
 \mu(s,z)\Longrightarrow z<s,
 \qquad \neg\mu(z,z).
 \tag{2}
\]
Define a noninjective map of variable labels by
\[
 f(4)=1,\qquad f(i)=i\quad(i\ne4).
 \tag{3}
\]
Translate a numerical formula recursively into ordinary classical
first-order arithmetic formulas:
\begin{align*}
 T_f(0)&=\top,\\
 T_f(M(a,b))&=\mu(x_{f(a)},x_{f(b)}),\\
 T_f(E(a,b))&=(x_{f(a)}=x_{f(b)}),\\
 T_f(U(a,B))&=\forall x_{f(a)}\,T_f(B),\\
 T_f(I(A,B))&=T_f(A)\to T_f(B).
\end{align*}
The last equation is understood semantically at the exceptional pair
$(A,B)=(0,0)$: both its sides are true. This suffices for soundness of
modus ponens. By (2), $T_f(1)=\mu(x_0,x_0)$ is false under every
assignment, even inside other binders. Thus the encoded negation,
conjunction and existential quantifier have their usual classical
semantics under this interpretation.

Every unchanged head (H1)--(H6) is valid under $T_f$. (H1) is
$(\exists y\,\forall y\,\psi)\to\psi$ over a nonempty domain;
(H2) is commutation of universal quantifiers, also valid if the two
labels become equal; (H3) and (H4) are ordinary substitution of equals
in the two argument positions of $\mu$; (H5) is the equality witness
$\exists y\,(z=y)$, or $\exists y\,(y=y)$ if the labels coincide;
and (H6) is universal distribution. The dominated propositional head
is a classical tautology for arbitrary values of all its formula
arguments, so its numerical restriction causes no problem. Truth is
valid. Modus ponens and generalization preserve validity.

It remains to check the two axiom families that mention numerical
freshness. Let $\operatorname{supp}(B)$ contain every variable label in
the decoded formula, including bound labels. The usual code induction
gives
\[
 v\in\operatorname{supp}(B)\Longrightarrow v<B.
 \tag{4}
\]
At $B=0$ the support is empty. The two exceptional codes in (3) have
\[
 1=M(0,0),\qquad 4=I(1,0),\qquad
 \operatorname{supp}(1)=\operatorname{supp}(4)=\{0\}.
 \tag{5}
\]
If $B\notin\{1,4\}$, then $f(B)=B$ and $f(v)\le v<B$ for every
$v\in\operatorname{supp}(B)$. If $B=1$ or $4$, equations (3)--(5)
give $f(B)=1\ne0=f(0)$. Hence in every case the variable
$x_{f(B)}$ is absent from $T_f(B)$. Every diagonal axiom
$B\to U(B,B)$ is therefore valid. This explicitly checks the codes
where the two labels were identified; injectivity of $f$ is unnecessary.

Now take an arbitrary padded gate and put $r=F(a,b,c)$.
By (1), $r\ge544$, so $f(r)=r$. Also $f(a),f(b)<r$.
By (4), the variable $x_r$ is absent from the translated predicate
$\psi=T_f(r)$. Thus its separator binder is genuinely fresh under
this interpretation. We verify the separation portion, including the
case $a\ne b$ but $f(a)=f(b)$.

Set $u=f(a)$ and $v=f(b)$, and fix an assignment of all variables
other than $x_r$. The membership-vacuity clause $T_f(V)$ is true
because $r\ne u,v$. Since this clause is true, the gate is equivalent
to
\[
 \exists x_r\bigl(T_f(S)\wedge T_f(P)\bigr).
 \tag{6}
\]
If $u\ne v$, write $A$ and $B$ for the current values of $x_u$
and $x_v$, and let $\psi(z)$ be the truth value of the translated
predicate after setting $x_v=z$, leaving the other parameters fixed.
Choose a finite-set code $s$ for
\[
 \{z:(\mu(A,z)\lor z=B)\wedge\psi(z)\}.
 \tag{7}
\]
This is a finite set: its elements are in the set coded by $A$ or
are equal to $B$. Assign $x_r=s$. Then, for every $z$,
\[
 \mu(A,z)\Longrightarrow(\psi(z)\Longrightarrow\mu(s,z)),
 \qquad \mu(s,z)\Longrightarrow\psi(z),
\]
which is exactly $T_f(S)$. The distinguished point $B$ in (7) gives
$\psi(B)\to\mu(s,B)$, which is $T_f(P)$.

If $u=v$, changing the bound variable in $S$ changes both original
occurrences together. Its first antecedent becomes $\mu(z,z)$ and
is always false by (2). Let $B$ be the current value of this shared
variable. Choose $s=2^B$ if $\psi(B)$ is true, and $s=0$ otherwise.
Then $\mu(s,z)\to\psi(z)$ holds for every $z$, while
$\psi(B)\to\mu(s,B)$ holds at the distinguished point. This again
proves (6). The predicate has no occurrence of $x_r$, so choosing
$s$ does not change these truth values. Every padded gate is valid.

We have established soundness of the whole modified calculus under
$T_f$. But $T_f(V(1,0,4))$ is
\[
 \mu(x_1,x_0)\to\forall x_1\,\mu(x_1,x_0).
\]
At $x_1=1,x_0=0$ its antecedent is true and its consequent is false
(take $x_1=0$ in the quantifier). Thus the old vacuity instance is
not valid. At the same assignment $V$ is false, so $V\to Q$ is true
regardless of $Q$ and $N(V\to Q)$ is false. Consequently
$T_f(G(1,0,4))$ is false as well. By soundness, neither formula is
derivable in $\mathcal C_{\mathrm{pad}}$.
\end{proof}

\subsection{Scope and consequences}

The failure is stronger than a missing proof of an obvious alpha-renaming
lemma: the particular old-gate recovery lemma is false for the stated
calculus. The padding permits the interpretation to identify two
numerically distinct variable labels that the original arbitrary-label
vacuity schema distinguishes. It is therefore circular to invoke full
first-order renaming in order to justify this replacement, unless that
renaming principle has been independently recovered from other axioms.

The theorem does not assert inconsistency of the padded calculus, nor
that it cannot be equiconsistent with PA by a different interpretation.
It does not rule out a different guardless redesign with an additional
independence or renaming mechanism. It rules out the claimed derivation
of all original gate instances from exactly the retained heads and the
padded gate family above.

A related earlier padding proposal in the corpus used
$K(D,H)=I(D+1,I(H,D+1))$ and $A^\star=I(K(D,H),A)$; it explicitly left
instance recovery open. The present theorem gives an explicit semantic
obstruction to the proposed recovery for the membership-tautology
wrapper, rather than merely noting a formalization gap. The supplied
\texttt{check\_codes.py} checks (5), $I(1,1)=16$, and
$F(0,0,0)=544$ exactly. These finite numeral checks supplement the
universal written proof and are not used as a substitute for it.

One possible further direction is a translation of PA proofs into a
restricted public image that never requires the missing raw gate. The
ConciseIPRA project demonstrates why completeness on a translated image
and completeness on arbitrary raw syntax must be distinguished. No such
translation for $\mathcal C_{\mathrm{pad}}$ is constructed here. This is
an open research direction, not a positive calibration theorem or a
claim that padding can be implemented below 194 states.
